\documentclass{article}
\usepackage{amsmath,amssymb}

\begin{document}

\begin{enumerate}
    \item Let \( f \colon \mathbf{R} \to \mathbf{R} \) be the function defined by \( f(x) = 3x+1 \).
    \begin{enumerate}
        \item Compute \( f^{-1}(\{1,5,8\}) \) (do not give a proof). \\
        To find \( f^{-1}(\{y\}) \) for each \( y \in \{1,5,8\} \), solve \( 3x + 1 = y \) for \( x \):
        \[
        f^{-1}(1) = \left\{ \frac{1-1}{3} \right\} = \{0\}, \quad
        f^{-1}(5) = \left\{ \frac{5-1}{3} \right\} = \left\{\frac{4}{3}\right\}, \quad
        f^{-1}(8) = \left\{ \frac{8-1}{3} \right\} = \left\{\frac{7}{3}\right\}.
        \]
        Thus, 
        \[
        f^{-1}(\{1,5,8\}) = \{0, \frac{4}{3}, \frac{7}{3}\}.
        \]

        \item Compute \( f^{-1}(W) \), where \( W = (4,\infty) \), and give a proof that your answer is correct. \\
        To find \( f^{-1}(W) \), solve \( 3x + 1 > 4 \) for \( x \):
        \[
        3x + 1 > 4 \Rightarrow 3x > 3 \Rightarrow x > 1.
        \]
        Therefore, \( f^{-1}(W) = (1, \infty) \). \\
        \textbf{Proof:} If \( x > 1 \), then \( 3x + 1 > 4 \), so every \( x \) in \( (1, \infty) \) maps to \( W \). Conversely, if \( y \in W \), then \( y > 4 \), and solving \( 3x + 1 = y \) gives \( x > 1 \).

        \item Compute \( f^{-1}(\mathbf{E}) \), where \( \mathbf{E} \) is the set of even integers, and give a proof that your answer is correct. \\
        Solving \( 3x + 1 = 2k \) for \( x \), where \( k \) is an integer:
        \[
        3x = 2k - 1 \Rightarrow x = \frac{2k-1}{3}.
        \]
        Since \( 2k-1 \) is always odd, \( x \) must be of the form \( \frac{\text{odd}}{3} \). Therefore, \( f^{-1}(\mathbf{E}) = \left\{ \frac{2k-1}{3} : k \in \mathbf{Z}, \ 3 \mid (2k-1) \right\}. \\
        \textbf{Proof:} Let $k$ be an integer such that $3$ divides \( 2k-1 \). Then \( x = \frac{2k-1}{3} \) is an integer. Hence, \( 3x+1 = 2k \), which is even, so every \( x \) in this set maps to \( \mathbf{E} \). Conversely, any even integer \( y = 2k \) can be written as \( 3x + 1 \) with \( x = \frac{2k-1}{3} \), provided \( 3 \) divides \( 2k-1 \).

    \end{enumerate}

    \item Let \( f \colon \mathbf{Z} \to \mathbf{Z} \) be the function defined by
    \[
    f(n) =
    \begin{cases}
    \frac{n}{2}, & \text{if } n \text{ is even} \\
    2n+4, & \text{if } n \text{ is odd}.
    \end{cases}
    \]
    Compute \( f^{-1}(\mathbf{E}) \). Prove that your answer is correct. \\


    We claim that \( f^{-1}(\mathbf{E})  = \mathbb{O} \cup 4\mathbb{Z}\).
    To find \( f^{-1}(\mathbf{E}) \), note that \( f(n) \) is even if \( n \) is even, since \( \frac{n}{2} \) is an integer. When \( n \) is odd, \( 2n+4 \) is always even. Thus, \( f^{-1}(\mathbf{E}) = \mathbf{Z} \). \\
    \textbf{Proof:} For any \( n \) in \( \mathbf{Z} \), whether \( n \) is even or odd, \( f(n) \) yields an even number. Hence, \( f^{-1}(\mathbf{E}) \) includes all integers.

    \item Let \( A \) and \( B \) be sets, and let \( X \) and \( Y \) be subsets of \( B \). Let \( f \colon A \to B \) be a function. Prove or disprove the following:
    \begin{enumerate}
        \item \( f^{-1}(X \cap Y) \subseteq f^{-1}(X) \cap f^{-1}(Y) \). \\
        \textbf{Proof:} If \( a \in f^{-1}(X \cap Y) \), then \( f(a) \in X \cap Y \), so \( f(a) \in X \) and \( f(a) \in Y \). Therefore, \( a \in f^{-1}(X) \) and \( a \in f^{-1}(Y) \), which implies \( a \in f^{-1}(X) \cap f^{-1}(Y) \).
        
        \item \( f^{-1}(X \cap Y) \supseteq f^{-1}(X) \cap f^{-1}(Y) \). \\
        \textbf{Proof:} If \( a \in f^{-1}(X) \cap f^{-1}(Y) \), then \( f(a) \in X \) and \( f(a) \in Y \), so \( f(a) \in X \cap Y \). Therefore, \( a \in f^{-1}(X \cap Y) \).
    \end{enumerate}
\end{enumerate}

\end{document}

